8.4 Real-valued Functions of Several Variables
477
    \begin{enumerate}
        \setcounter{enumi}{5}
        \renewcommand{\theenumi}{\alph{enumi}}
        \item Let
        \[
        D^\mathbf{\alpha} f(\mathbf{x}) := \frac{\partial^{|\mathbf{\alpha}|} f}{(\partial x^1)^{\alpha_1} \ldots (\partial x^m)^{\alpha_m}}(\mathbf{x}).
        \]
        Show that if $f \in C^{(k)}(G; \mathbb{R})$, then the equality
        \[
        \sum_{i_1 + \ldots + i_m = k} \frac{k!}{i_1! \ldots i_m!} f^{(i_1, \ldots, i_m)} h^{i_1} \ldots h^{i_m} = \sum_{|\mathbf{\alpha}|=k} \frac{k!}{\mathbf{\alpha}!} D^\mathbf{\alpha} f(\mathbf{x}) \mathbf{h}^\mathbf{\alpha},
        \]
        where $\mathbf{h} = (h^1, \ldots, h^m)$, holds at any point $\mathbf{x} \in G$.
        \item Verify that in multi-index notation Taylor's theorem with the Lagrange form
        of the remainder, for example, can be written as
        \[
        f(\mathbf{x} + \mathbf{h}) = \sum_{|\mathbf{\alpha}|=0}^{n-1} \frac{1}{\mathbf{\alpha}!} D^\mathbf{\alpha} f(\mathbf{x}) \mathbf{h}^\mathbf{\alpha} + \sum_{|\mathbf{\alpha}|=n} \frac{1}{\mathbf{\alpha}!} D^\mathbf{\alpha} f(\mathbf{x} + \theta \mathbf{h}) \mathbf{h}^\mathbf{\alpha}.
        \]
        \item Write Taylor's formula in multi-index notation with the integral form of the
        remainder (Theorem 4).
    \end{enumerate}
    \item Prove the following generalization of Rolle's theorem for functions of several
    variables.
    If the function $f$ is continuous in a closed ball $B(\mathbf{0}; r)$, equal to zero on the
    boundary of the ball, and differentiable in the open ball $B(\mathbf{0}; r)$, then at least one of
    the points of the open ball is a critical point of the function.
    \item Verify that the function
    \[
    f(x, y) = (y - x^2)(y - 3x^2)
    \]
    does not have an extremum at the origin, even though its restriction to each line
    passing through the origin has a strict local minimum at that point.
    \item The method of least squares. This is one of the commonest methods of processing
    the results of observations. It consists of the following. Suppose it is known that the
    physical quantities $x$ and $y$ are linearly related:
    \[
    y = ax + b \tag{8.79}
    \]
    or suppose an empirical formula of this type has been constructed on the basis of
    experimental data.
    Let us assume that $n$ observations have been made, in each of which both $x$ and $y$
    were measured, resulting in $n$ pairs of values $x_1, y_1; \ldots; x_n, y_n$. Since the measure-
    ments have errors, even if the relation (8.79) is exact, the equalities
    \[
    y_k = ax_k + b
    \]
    may fail to hold for some of the values of $k \in \{1, \ldots, n\}$, no matter what the coeffi-
    cients $a$ and $b$ are.
    The problem is to determine the unknown coefficients $a$ and $b$ in a reasonable
    way from these observational results.
    Basing his argument on analysis of the probability distribution of the magnitude
    of observational errors, Gauss established that the most probable values for the co-
    efficients $a$ and $b$ with a given set of observational results should be sought by use
    of the following least-squares principle:
    If $\delta_k = (ax_k + b) - y_k$ is the discrepancy in the $k$th observation, then $a$ and $b$
    should be chosen so that the quantity
    \[
    \Delta = \sum_{k=1}^n \delta_k^2
    \]
    that is, the sum of the squares of the discrepancies, has a minimum.
    \begin{enumerate}
        \item[(a)] Show that the least-squares principle for relation (8.79) leads to the following
        system of linear equations
        \[
        \begin{cases}
        [x_k, x_k]a + [x_k, 1]b = [x_k, y_k], \\
        [1, x_k]a + [1, 1]b = [1, y_k],
        \end{cases}
        \]
    \end{enumerate}